\begin{abstract}
An ordered assignment of independently shifted finite-difference weights
contains information that its determinantal symmetrization erases. We
recover that information explicitly. For two particles, two finite
summations by parts reduce the ordered double sum to a single two-shift
zeta kernel. This gives its joint meromorphic continuation, every polar
coefficient, all convergent positive-integer-order evaluations and their
finite parts, a divided-difference
digamma generating function, and a Gamma primitive. Oblique poles explain
direction-dependent values at parameters where every pointwise weight
vanishes. A cutoff generating function is then proved on an entire
two-order parameter domain: its sole depth-two polylogarithm has the
factor $(1-z)^2$, which explains the endpoint reduction. Beta
parametrization yields ordinary convergent identities involving every
generalized Stieltjes constant. At arbitrary particle number, transferring
the stencils to the ordered Vandermonde square localizes the sum to
collision blocks; an explicit finite algorithm uses at most
$\lfloor(r+1)/2\rfloor$ free summation indices. A separate Gamma-integral
basis solves the fixed-gap null equation for unequal even harmonic
tails. Its exact normalization and nonzero determinant asymptotic prove
eventual optimal reconstruction on every fixed-gap ray; the first
relative correction is an explicit harmonic-number expression.
All formulas retain their convergence domains, confluence conventions,
and finite-part coordinates. Exact algebraic tests and numerical checks
accompany the proofs.
\end{abstract}

\section{The source questions and the mathematical setting}
\label{sec:scope}

\subsection{A continuation of the current material}

The source snapshot is the ProveIt revision \repo, obtained on
11 October 2026. Both requested directories were inspected: the canonical
manuscript and the seven then-present incoming archives. The manuscript's
claim ledger remains the authority for accepted canonical results
\cite{repo:manuscript,repo:ledger}; an incoming report's delivery status
does not by itself amount to canonical proof acceptance.

The most direct target is the incoming \emph{Finite-Difference
Determinants} report \cite{repo:finite}. Its section entitled
``Unsymmetrized mixed-order sums'' asks which additional coordinates are
needed for an individual ordering of independent spectral labels.
Its existing polarization theorem gives only the sum over assignments.
It also asks for a finite description of noninteger-order cutoff
generators and for extensions to other annihilating stencils. These
questions are mathematically related: shifting the finite differences
onto the ordered summation domain exposes the boundaries at which
the extra coordinates live.

The present two-particle result answers the individual-ordering question
on an open complex domain, not only at isolated integer values.
It shows that a one-index kernel is enough. That kernel has explicit
Hurwitz expansions and, at positive integer orders, a finite polygamma
evaluation. The all-particle result is different: it gives a finite
localization formula and a bound on the number of free indices, rather
than incorrectly extrapolating a universal single-zeta formula.

A second target comes from the harmonic-tail Gamma continuation
\cite{repo-gamma}. Its exact fixed-gap reduction leaves an
$m$-dimensional null equation, where the two tail orders differ by
$2m$. Section~\ref{gb:section} constructs a basis by ordinary integrals
against $|\Gamma(\tfrac14+\tfrac{ix}{2})|^{4m}$. The distinct
decay scales of these integrals produce a nonzero limiting Vandermonde
determinant. This proves that, for every fixed $m$, the first $2b$
moments form the optimal formal prefix for all sufficiently large
$(p,q)=(2b-2m,2b)$. The theorem is complementary to the existing
eventual result with the smaller order fixed \cite{repo:quartic}. Neither assertion alone
settles every unequal pair.

The current sources also make clear why old open labels cannot simply
be reused. Several finite-order harmonic-tail reconstruction cases have
already received proofs in recent reports. The short Gaussian $S_6$
and revised $S_8$ evaluations, by contrast, remain conjectural in the
current canonical manuscript \cite{repo:ledger,repo:gaussian}. None of
the ordered identities below is represented as a proof of either
Gaussian evaluation.

\subsection{Classical ingredients and the scope of originality}

Finite summation by parts, finite differences, partial fractions,
Hurwitz continuation, and Gamma differentiation are classical.
The general rational-summability role of discrete residues is treated
by Chen--Singer and Arreche--Sitaula \cite{chen-singer,arreche-sitaula}.
Multivariable zeta functions can have intersecting polar hyperplanes
and direction-sensitive Laurent behavior; this phenomenon is not new
in itself \cite{matsumoto-ono-wakabayashi}. Beta parametrization is also
an established method for zeta products \cite{shpot-paris}.

The contributions claimed here are the explicit ordered boundary
identity, its complete consequences in this family, the finite cutoff
polylogarithm formula and its endpoint cancellation, and the
all-particle collision-block reduction. Their derivations were checked
against the inspected sources. The literature search does not establish
worldwide historical priority. In particular, finite polygamma
reduction of a rational one-index summand is a classical step; it is
the preceding reduction of the ordered ensemble that supplies the
new application.

The Gamma basis, its normalization, and its fixed-gap consequences form
the second group of contributions. Classical Fourier transforms of
Gamma products provide useful context \cite{ramanujan-transform};
the result here concerns a specific higher-power weight, its
sinh/tanh modes, and the repository's exact reconstruction determinant.

\subsection{Notation and analytic conventions}

For $\Rea a>0$ we use the principal logarithm in
$(n+a)^{-s}=\exp(-s\log(n+a))$. Thus every shifted denominator in
the original series stays in the right half-plane. Write
\[
 (s)_j=s(s+1)\cdots(s+j-1),\qquad (s)_0=1,
 \qquad \psi(a)=\frac{\Gamma'(a)}{\Gamma(a)}.
\]
For positive integers $p$,
\[
 H_n^{(p)}=\sum_{k=1}^{n}k^{-p},\qquad H_0^{(p)}=0.
\]
The same finite sum defines $H_n^{(s)}$ for complex $s$.
We use the Stieltjes convention
\begin{equation}\label{eq:stieltjes-convention}
 \zeta(1+u,a)=\frac1u+
    \sum_{m=0}^{\infty}\frac{(-1)^m\gamma_m(a)}{m!}u^m,
 \qquad \gamma_0(a)=-\psi(a).
\end{equation}
Derivatives of $\zeta$ without an argument qualifier differentiate its
first variable. Standard relations used below include
\begin{equation}\label{eq:classical-ladders}
 \partial_a\zeta(s,a)=-s\zeta(s+1,a),\qquad
 \zeta(p,a)=\frac{(-1)^p}{(p-1)!}\psi^{(p-1)}(a)\quad(p\ge2),
\end{equation}
and $\zeta'(0,a)=\log\Gamma(a)-\tfrac12\log(2\pi)$
\cite{dlmf:hurwitz,dlmf:gamma}.

An equality continued from a normally convergent domain is always an
equality of those continuations. It need not equal the result of
substituting exceptional parameters into every summand first.
Likewise, a finite part is the constant in a declared Laurent
coordinate; a hard cutoff is a different operation unless their
agreement is proved.

For $|z|<1$ our multiple-polylogarithm convention is
\begin{equation}\label{eq:mpl-convention}
 \Li_{t,s}(z,1)=\sum_{k>j\ge1}\frac{z^k}{k^t j^s}.
\end{equation}
The exponential factor makes this entire in the two order variables.
This is independent of the convergence restrictions of its value at
$z=1$.
